\section{n4. Positive envelopes at both ends of a noninjective counterexample}
\paragraph{Supports.} G331, G345, n1 and n2.
For EVERY field $k$, take $R=k[x,y]_{(x,y)}$, $Q=R/(xy)$, gauge $\mathfrak m=(x,y)$. $R$ is excellent normal regular of dimension two, $Q$ excellent reduced of dimension one. Its finite normalization has two DVR centres with residue degree one. The FULL $X_Q$ has exactly
\[
 w_x(x)=1,w_x(y)=\infty;\qquad w_y(x)=\infty,w_y(y)=1.
\]
Restriction is injective on these points, with closed two-point image in $X_R$. It is not onto: $o=\operatorname{ord}_{\mathfrak m}$ is a further point with support zero. The graded-domain argument makes $o$ a valuation.

On ALL $X_R$ let $f=-1/v(x)$, with reciprocal at infinity zero. Since $v(x)\ge1$, this is continuous bounded in $[-1,0]$. It has values $-1$ at $o,w_x$ and zero at $w_y$. As in (3), $J\subset\mathfrak m^{o(J)}$ gives $u(v)\le u(o)$ for EVERY $C_R,D_R$ potential on EVERY support. Thus an obstacle competitor is at most $-1$ everywhere. The genuine primary profile $-v(\mathfrak m)=-1$ attains it:
\[
 P_R^Cf=P_R^Df=-1\quad\hbox{on FULL }X_R. \tag{5}
\]
On $Q$, the genuine primary ideals $J_N=(x^N,y)Q$ give profiles $-w(J_N)/N=(-1,-1/N)$, uniformly converging on the full two-point link to $h=(-1,0)$. G331 comparison gives $h\in C_Q\subset D_Q$. It is the obstacle itself, hence
\[
 P_Q^C(f\rho)=P_Q^D(f\rho)=(-1,0). \tag{6}
\]
Equations (5)--(6) have positive continuous geometric envelopes on BOTH ends and an exact gap one at $w_y$. This failure is not an envelope-existence gap.

There is no $D_R$ extension of $h$: a value zero at $\rho w_y$ forces identically zero by G331's normal zero dichotomy, contradicting the value at $\rho w_x$. Fibre constancy on an injective boundary map is therefore insufficient for positive extension when the ring map is noninjective. Factoring through the exact image $Q$ makes the residual injection an identity; it gives no extension to $R$.

Each ideal $J_N$ lifts to the primary $R$ ideal $(x^N,y)$, but its ambient profile $g_N=-\min(v(x),v(y)/N)$ violates $f$ at $o$ for $N>1$. The $g_N$ increase pointwise to zero where $v(y)<\infty$, and to $-1$ where $v(y)=\infty$. Monomial valuations $(v(x),v(y))=(1,M)$ tend to $\rho w_x$; the limit function is zero along them and $-1$ at the endpoint. It is not USC, hence not $D_R$. Increasing closure cannot repair extension.

For a primary $\mathfrak b$, transport by $T(v)=v/v(\mathfrak b)$ and divide functions and obstacles by the same positive factor. This is an order bijection of competitors, so transports (5)--(6). In homogeneous notation $f_{\mathfrak b}(w)=-w(\mathfrak m)^2/w(x)$ and the ambient envelope is $-w(\mathfrak m)$. At the transported $y$ branch the gap is $1/\operatorname{ord}_y(\mathfrak bQ)>0$, one function factor.

G345 gives $\mu_R(-1)=\delta_o$ by degree one of the exceptional $\mathbb P^1$, but $\mu_Q()=\delta_{w_x}+\delta_{w_y}$ by zero-arity DVR Cartier degrees. The former assigns zero to the closed image; the latter has mass two. These have different arities. The quotient has zero fixed ambient degree-two contribution, while its intrinsic degree-one reference is nonzero. This is consistent with n5, not a failure of fixed-ambient projection.
